\paragraph{E.3n (Proved: a torsion-valued abelian-family map to $P$ over
$\mathbb C$ is zero).}
Let $X/\mathbb C$ be smooth projective, let $A$ be an abelian
variety, and let $Z\in\mathrm{CH}^i(A\times X)$ be normalized.
Then
\[
 u_Z\otimes\mathbb Q=0\quad\Longrightarrow\quad u_Z=0.
\]
The hypothesis means that every value of $u_Z$ is torsion;
it does not assume a uniform exponent. We prove the required
uniformity from finite-type parameter spaces for the defining
Albanese relations. Uncountability of $\mathbb C$ is essential
to this proof, as the example E.3f over $\overline{\mathbb F}_p$
also shows.

\emph{Moving the points before parameterizing the relations.}
For a smooth projective connected $T$ and a nonempty open
$T^\circ\subset T$, every class in $\mathrm{CH}_0(T)$ has a
zero-cycle representative supported in $T^\circ$; equivalently,
the map $Z_0(T^\circ)\to\mathrm{CH}_0(T)$ from the group of
cycles is surjective. We are not using a covariant Chow-group
map for an open immersion.
Indeed, for each excluded point choose an integral curve
through it whose generic point lies in $T^\circ$, normalize
the curve, and use Riemann--Roch to express that point as a
difference of effective divisors avoiding the finitely many
points outside $T^\circ$. Pushing to $T$ proves the assertion.
For dimension zero it is immediate.

Represent a correspondence $W\in\mathrm{CH}^i(T\times X)$
by a finite signed sum of integral subvarieties. Those which
do not dominate $T$ act trivially on $\mathrm{CH}_0(T)$:
move the zero-cycle away from their images. Each dominating
component is flat over a common nonempty open $T^\circ$
after generic flatness. We can therefore replace any
Albanese-kernel zero-cycle $\nu$ by a rationally equivalent
zero-cycle supported in $T^\circ$. Its Albanese class remains
zero. Now $W_*\nu$ is represented by the actual signed
sum of flat fibres $W_t$ at its finitely many support points.
Thus every generator of $N^i(X)$ has this form; we need not
parameterize arbitrary refined intersections at bad fibres.

\emph{Countably many finite-type spaces parameterize these generators.}
Embed $X$ once in projective space. We describe all data as
the projective embedding dimension, Hilbert polynomials,
finite lengths, and integer coefficients vary. There are
only countably many such choices.
Throughout we replace the parameter bases by their reductions
and use finite locally closed stratifications when required
for Hilbert--Chow, constant fibre dimension, or multiplicity.
This does not lose any complex-point datum; allowing all the
strata preserves countability and constructibility.

First use Hilbert schemes to parameterize closed embedded
smooth projective geometrically connected varieties $T$.
The locus with these fibre properties is open in the
appropriate Hilbert scheme; allowing the finitely many
components and locally closed strata gives finite-type bases.
Pass to the universal $T$ itself as a base to mark an origin
$t_0$. The pulled-back family has a section. Its relative
Albanese is an abelian scheme, with its pointed map, compatible
with base change. The precise characteristic-zero existence
and base-change statement is
\citet[Theorem 2.2]{cyJavanpeykarMathur2022}; the marked section
turns their Albanese torsor into the pointed group scheme.

Relative Hilbert schemes for $T\times X$ parameterize the
finitely many components of $W$. Restrict to the locus where
the marked evaluation points lie in the fibre-flat locus.
This is an algebraic constructible condition. More explicitly,
the universal subschemes are flat over their Hilbert base;
the local criterion for flatness identifies fibrewise flatness
over $T$ with flatness of their projections over the universal
$T$, and the bad locus has closed image because these
projections are projective. Finite stratification also fixes
the dimensions and multiplicities of the resulting fibre
cycles. Mark the finitely many points $t_1,\ldots,t_s$ by
the relative power $T^s$. For fixed integer weights $e_j$
with $\sum_j e_j=0$, the condition
\[
            \sum_j e_j\operatorname{alb}_{T}(t_j)=0
\]
is closed: it is the inverse image of the zero section under
an algebraic morphism to the relative Albanese. On this
finite-type locus, the signed flat-fibre cycle
$\sum_j e_jW_{t_j}$ is an algebraic family of cycles on $X$.
Its positive and negative parts are represented, after finite
stratification, by morphisms to fixed-degree Chow schemes.
Equivalently one can apply the Hilbert--Chow morphism to each
flat family and then use the algebraic addition of cycles.
Allowing finite tuples of these data parameterizes every
finite sum of generators of $N^i(X)$.

\emph{Rational-equivalence witnesses also have countable finite-type data.}
An integral rational equivalence is a finite sum of signed
differences $[V_0]-[V_\infty]$, where
$V\subset\mathbb P^1\times X$ is an integral closed subscheme
dominating $\mathbb P^1$. Such a subscheme is flat over the
smooth curve $\mathbb P^1$. The equivalence of this description
with principal-divisor rational equivalence follows by closing
the graphs of the finitely many rational functions and pushing
forward; it is proved in
\citet[Lemma 43.9.1, Tag 0AZI]{cyStacksRationalEquivalence}.
Hilbert schemes of $\mathbb P^1\times X$, indexed by their
countably many Hilbert polynomials, therefore parameterize
all witnesses. We restrict to the locus of flat projection
to $\mathbb P^1$, evaluate at $0,\infty$, and apply
Hilbert--Chow and cycle addition. Equality of the two resulting
effective sums is a closed condition in a fixed Chow scheme.
Finite tuple lengths and integer signs give only countably
many more choices. Consequently the locus that two
specified algebraic cycle families are rationally equivalent
is a countable union of constructible sets, by Chevalley's
theorem. Applying this to the preceding finite sums proves
that membership in $N^i(X)$, in any fixed algebraic family
of signed cycles, is countably constructible.

For the same Hilbert--Chow construction of countable
rational-equivalence witnesses, with its maps written out,
see the opened primary proof of
\citet[Lemma 4.4, pp.~311--313]{cySchoen1993}.
Our additional parameter spaces above impose the relative
Albanese-kernel condition; that condition is not asserted
to be part of Schoen's lemma.

All varieties and cycles in a particular relation occur in
this enumeration: projective embeddings and Hilbert polynomials
are finite data, the Albanese condition is exact, the moving
step preserves the given Chow class integrally, and every
rational equivalence uses finitely many functions. Thus the
construction does not replace $N$ by its saturation, closure,
or a subgroup defined only by regular invariants.

\emph{Apply uncountability and the group law.}
Choose a nonempty Zariski open $U\subset A$ where the finitely
many dominating components of a representative of $Z$ are
flat; vertical components vanish on the fibres there. Fix
any cycle representative of its refined fibre at the origin.
Then $Z_a-Z_0$ is a signed algebraic cycle family for $a\in U$.
For every $m>0$, the preceding construction gives
\[
 \{a\in U(\mathbb C):m u_Z(a)=0\}
                 =\bigcup_{j\ge1}C_{m,j}(\mathbb C),
\]
where every $C_{m,j}$ is constructible in $U$.
If every $u_Z(a)$ is torsion, the countable collection of all
$C_{m,j}$ covers $U(\mathbb C)$. A complex irreducible variety
cannot be the union of countably many proper closed algebraic
subsets: this follows, for example, from the Baire theorem on
a coordinate ball. Therefore some $C_{m,j}$ is Zariski dense
and contains a nonempty open $V\subset A$.

The full kernel of $m u_Z$ is a subgroup of $A(\mathbb C)$,
so it contains $V(\mathbb C)-V(\mathbb C)=A(\mathbb C)$.
The equality follows because $V$ meets every translate of
$V$ in the irreducible variety $A$. Thus $m u_Z=0$ with one
fixed positive integer $m$. Divisibility of $A(\mathbb C)$
then gives $u_Z=0$, as claimed.

Combining this result with E.3c and E.3m yields an exact
further reduction: for fixed $X,i$, the integral family
converse, the rational family converse, the rational
ordinary Abel--Jacobi injectivity on $P^i(X)\otimes\mathbb Q$,
and the integral equality $N^i(X)=\ker\psi_W$ are equivalent.
Their universal truth is refuted by the specified
codimension-four family of E.3u. Ordinary integral Abel--Jacobi
injectivity is a different assertion, refuted in E.3h--k.
